\begin{pastexam}[label={px:gw3-097-ds-05}]{關務三等資料處理 97 年 第五題}
（申論；15 分）

以插入法處理下列 10 個數字，排出由小到大的順序，寫出其演算法。（15 分）

數字為：100，31，60，5，23，89，26，1，50，55
\end{pastexam}

\begin{pastsoln}
\textbf{演算法（插入排序）。} 把陣列分成左邊「已排好」和右邊「還沒處理」兩段，一開始左段只有第 $1$
個數。 每次取右段的第一個數 $\mathit{key}$，從左段的尾端往前比，比 $\mathit{key}$ 大的都往右挪一格，
直到遇到不比它大的數，把 $\mathit{key}$ 放進空出來的位置（資結書虛擬碼~1.3）。
\begin{pseudo}{插入排序}
\Procedure{Insertion-Sort}{$A, n$}
    \For{$i \gets 2$ \To $n$}
        \State $\mathit{key} \gets A[i]$
        \State $j \gets i - 1$
        \While{$j \ge 1$ \AND $A[j] > \mathit{key}$}
            \State $A[j + 1] \gets A[j]$ \Comment{比 $\mathit{key}$ 大的往右挪}
            \State $j \gets j - 1$
        \EndWhile
        \State $A[j + 1] \gets \mathit{key}$ \Comment{放進空位}
    \EndFor
\EndProcedure
\end{pseudo}

\textbf{處理過程}（每一列是插入一個數之後的陣列，粗體是已排好的部分，框起來的是剛插入的數）：
\begin{center}
\small
\setlength{\tabcolsep}{4pt}
\begin{tabular}{ccl}
\toprule
$i$ & 插入 & 之後的陣列 \\
\midrule
初始 & & \textbf{100}, 31, 60, 5, 23, 89, 26, 1, 50, 55 \\
2 & 31 & \textbf{\fbox{31}, 100}, 60, 5, 23, 89, 26, 1, 50, 55 \\
3 & 60 & \textbf{31, \fbox{60}, 100}, 5, 23, 89, 26, 1, 50, 55 \\
4 & 5 & \textbf{\fbox{5}, 31, 60, 100}, 23, 89, 26, 1, 50, 55 \\
5 & 23 & \textbf{5, \fbox{23}, 31, 60, 100}, 89, 26, 1, 50, 55 \\
6 & 89 & \textbf{5, 23, 31, 60, \fbox{89}, 100}, 26, 1, 50, 55 \\
7 & 26 & \textbf{5, 23, \fbox{26}, 31, 60, 89, 100}, 1, 50, 55 \\
8 & 1 & \textbf{\fbox{1}, 5, 23, 26, 31, 60, 89, 100}, 50, 55 \\
9 & 50 & \textbf{1, 5, 23, 26, 31, \fbox{50}, 60, 89, 100}, 55 \\
10 & 55 & \textbf{1, 5, 23, 26, 31, 50, \fbox{55}, 60, 89, 100} \\
\bottomrule
\end{tabular}
\end{center}
說明其中幾步：$i = 4$ 插入 $5$ 時，$100, 60, 31$ 都比 $5$ 大，依序右移，$j$ 走到 $0$，$5$ 放到第 $1$
格。 $i = 6$ 插入 $89$，只比 $100$ 小，$100$ 右移一格，$89$ 放在它前面。 $i = 8$ 插入 $1$，前面 $7$
個都右移。 $i = 10$ 插入 $55$，$100, 89, 60$ 右移，遇到 $50 \le 55$ 停下。

結果：$1, 5, 23, 26, 31, 50, 55, 60, 89, 100$。 各步的比較次數依序是 $1, 2, 3, 4, 2, 5, 7, 4, 4$，
共 $32$ 次（含每次停下來的那一次比較；$j$ 走到 $0$ 時不用比）。 插入排序最壞 $O(n^2)$、最好（已排好）
$O(n)$，是穩定、原地的排序。
\end{pastsoln}
